% Pursuing Stacks (cotes 134-2 to 134-8): archive cote 134-5, site page 128 (PDF page 129), typescript page 315.
% typescript. Draft compared: ps5.tex:1130-1184.
% AI-transcribed by Claude Opus 5.5 (claude-opus-5-5, adaptive thinking, effort high), Message Batches API, batch ps3 (msgbatch_01XdFmuVVBAcgc7EFxdtCqsj), 2026-11-12; not yet checked by eye.
\scanpage{134-5/128}
\mspage{315}

$A^{\wedge}_{\mathrm{fp}}$ . Thus the present question about finiteness conditions, seems
to be related (possibly) to the previous one about existence of various
types of test functors.

The condition for a strict test category A we are lookihg at is
surely satisfied, besides \hand{$\Delta$}, by the cubical and hemispherical test
categories \hand{$\Box$} and \hand{$\unsure{\bigcirc}$}, and surely also by any finite products of these.
I add this comment, of course, in order to ``push through'' the point
that \inserted{not any more} with respect to finiteness conditions on homotopy types, \over{as}{than}
(presumably at least, for the time being) in \deleted{no}\inserted{any} other essential respect
concerning the ability \xout{of} for expressing basic situations and facts in
homotopy or cohomology theory, the category of simplices stands singled
out by itself from all other test categories. N\hand{o}r does it seem that
the ``trinity''
\[ \hand{$\Delta$}\ ,\ \hand{$\Box$}\ ,\ \hand{$\unsure{\bigcirc}$} \]
has this property, with the only exception so far, possibly, of the
Dold-Puppe theorem\over{.}{,}\hand{(} as no other test category except these is known to
me for which a Dold-Puppe theorem in it's strict form holds true
(compare reflections section 71)\hand{)}.

8) I could make the same point in favor of more general test cate-%
gories than the trinity above, when it comes to the existence of an
algorithm\xout{\unsure{s}} for computing homology and cohomology groups, using suitab-%
le \uline{boundary} \uline{operators}. What is meant by these is clear of course for
the three types of complexes, but then it extends in an obvious way to
multicomplexes too - which means that \xout{xx}\inserted[typed]{for} the test categories deduced
from the trinity by taking finite products, we still get an algorithm\xout{\unsure{s}}
for cohomology via boundary operators. Of course, for any test catego-%
ry A , using a test functor (if we can find one) from one of the
three above (say) into $A^{\wedge}$ will allow us to reduce ``computation'' of
homology and cohomology invariants in terms of a model in $A^{\wedge}$, to the
case of the corresponding type of complexes - hence again an algo-%
rithm\deleted{\unsure{s}} (similar to the familiar one of computing the cohomomology of an
object of (Cat) \xout{xx} semisimplicially, via the nerve\hand{)}. But this is
cheating of course ! The question I want to raise here is about
\inserted{existence of} \hand{``}boundary operations\hand{''} \xout{\unsure{from A xx}}\inserted{in} A, similar to the familiar ones used for
the three \inserted{basic} types of complexes, and allowing to compute the homology and
cohomology groups of an object of $A^{\wedge}$ in the usual way, involving sui-%
table \uline{signs} + or - associated to the various boundary operations.

\margin{128}
